\begin{pastexam}[label={px:ncu-107-dsa-01}]{中央資工所 107 學年度 第 1 題}
（單選；每題 4 分，答錯倒扣一分）

Consider the following recursive sorting algorithm to sort an array A[\,] of $n$ numbers:
\begin{quote}
E-SORT(A, i, j)\\
1.\quad \textbf{if} A[i] $>$ A[j]\\
2.\quad\quad \textbf{then} exchange A[i] $\leftrightarrow$ A[j];\\
3.\quad \textbf{if} i + 1 $\ge$ j\\
4.\quad\quad \textbf{then} return;\\
5.\quad k $\leftarrow \lfloor$(j - i + 1)/3$\rfloor$;\qquad // Round down.\\
6.\quad E-SORT(A, i, j - k);\quad // Sort first two-thirds.\\
7.\quad E-SORT(A, i + k, j);\quad // Sort last two-thirds.\\
8.\quad E-SORT(A, i, j - k);\quad // Sort first two-thirds again.
\end{quote}
What is the time complexity of the algorithm if it is called as: E-SORT(A, 1, $n$)?

(a)~$\Theta(n)$ \qquad (b)~$\Theta(n\log n)$ \qquad (c)~$\Theta(n^{\log_2 3})$ \qquad
(d)~$\Theta(n^2)$ \qquad (e)~$\Theta(n^{\log_{3/2} 3})$
\end{pastexam}

\begin{pastsoln}
答案是 (e)。

\textbf{寫出遞迴式}（資結書第~2.1~節）。 令 $T(n)$ 是處理 $n = j - i + 1$ 個元素的時間。
第 1--5 行是常數時間。 $k = \lfloor n/3 \rfloor$，第 6、8 行處理 $A[i..j-k]$，有
$n - \lfloor n/3 \rfloor = \lceil 2n/3 \rceil$ 個元素；第 7 行處理 $A[i+k..j]$，也是
$n - k = \lceil 2n/3 \rceil$ 個。 $n \le 2$ 時在第 4 行直接返回。 所以
\[
T(n) = 3\,T\bigl(\lceil 2n/3 \rceil\bigr) + \Theta(1), \qquad T(1) = T(2) = \Theta(1).
\]

\textbf{解遞迴式}。 用主定理（資結書定理~2.4），$a = 3$，$b = 3/2$，$f(n) = \Theta(1)$。
\[
n^{\log_b a} = n^{\log_{3/2} 3}, \qquad \log_{3/2} 3 = \frac{\ln 3}{\ln 1.5} \approx \frac{1.0986}{0.4055}
\approx 2.71 .
\]
$f(n) = \Theta(1) = O(n^{\log_{3/2} 3 - \varepsilon})$（取 $\varepsilon = 1$ 就行），是第一種情況，
\[
T(n) = \Theta\bigl(n^{\log_{3/2} 3}\bigr) \approx \Theta(n^{2.71}).
\]
天花板不影響結果（資結書第~2.1.6~節）。 也可以用遞迴樹看：第 $d$ 層有 $3^d$ 個呼叫，每個
$\Theta(1)$；樹高 $\log_{3/2} n$，葉子有 $3^{\log_{3/2} n} = n^{\log_{3/2} 3}$ 個，葉子那一層
就決定了總和。

逐項檢查：
\begin{itemize}
\item (a) $\Theta(n)$、(b) $\Theta(n \log n)$：太小。 光是遞迴呼叫的次數就有 $\Theta(n^{2.71})$ 次。
\item (c) $\Theta(n^{\log_2 3}) \approx \Theta(n^{1.585})$：這是 $T(n) = 3T(n/2) + \Theta(1)$ 的解，
把子問題當成一半了。 這裡每個子問題是三分之二，比一半大，錯。
\item (d) $\Theta(n^2)$：$n^{2.71}$ 比 $n^2$ 高一階，錯。
\item (e) $\Theta(n^{\log_{3/2} 3})$：對。
\end{itemize}
這個演算法就是 Stooge sort。 它排得對（先排好前 $2/3$，最大的 $\lceil n/3 \rceil$ 個數一定都在
後 $2/3$；排好後 $2/3$，它們就到了正確位置；最後再排一次前 $2/3$），但比插入排序的
$\Theta(n^2)$ 還慢。
\end{pastsoln}
